\section{Generic Unfaithfulness and a Minimal Example}

\subsection{A conditional structural-unfaithfulness result}

The faithfulness condition places a strong constraint on the reduction $\Pi$: whenever
distinct histories are identified at the reduced level, their future responses to all
admissible interventions must remain indistinguishable. This can hold only if the
discarded distinctions are dynamically irrelevant for the specified intervention class
$\mathcal{I}$.

Rather than appealing to a measure-theoretic notion of typicality, we state a conditional
\emph{structural diagnostic}: if a reduction discards degrees of freedom that can become
dynamically relevant under interventions that are standard in open-system practice
(including free evolution), then the reduction fails to be faithful. In this sense,
faithfulness requires special structural coincidences (e.g., absence of correlations)
that are unstable under perturbations of histories, interactions, or intervention
richness.

\begin{proposition}[Unfaithfulness from discarded correlations under free evolution]
\label{prop:unfaithful-from-correlations}
Assume a decomposition of the underlying system into retained and discarded degrees of
freedom (e.g., $S$ and $E$), and suppose:
\begin{enumerate}
  \item (\emph{Degenerate reduction}) There exist histories $h_1 \neq h_2$ such that
        $\Pi(h_1)=\Pi(h_2)=r$, but $h_1$ and $h_2$ differ in correlations between retained
        and discarded degrees of freedom.
  \item (\emph{Intervention richness}) The intervention class $\mathcal{I}$ contains at least
        one ``passive'' intervention $I_{\Delta t}$ corresponding to free evolution of the
        full retained+discarded system over a duration $\Delta t$ (no access to discarded
        degrees of freedom by the experimenter).
  \item (\emph{Correlation sensitivity}) For the above $h_1,h_2$, free evolution yields
        successor histories $h_1' , h_2'$ with distinct reduced outcomes:
        $\Pi(h_1') \neq \Pi(h_2')$.
\end{enumerate}
Then $\Pi$ is unfaithful with respect to $\mathcal{I}$.
\end{proposition}

\begin{remark}[Interpretive status]
This result is diagnostic rather than prohibitive: it identifies sufficient structural
conditions under which algebraically closed reduced descriptions cannot support
counterfactual prediction under the intervention class $\mathcal{I}$. It does not claim
that such descriptions must always be abandoned, nor that they fail outside these
conditions.
\end{remark}

\begin{remark}[Meaning of structural genericity]
Genericity here is understood \emph{structurally} rather than statistically:
it refers to the instability of faithfulness under perturbations to correlation
structure or intervention classes, not to the frequency with which approximate
$\Sigma_1$ descriptions succeed in controlled experimental regimes.
\end{remark}

\begin{remark}[Regimes of $\Sigma_1$ adequacy]
Faithful $\Sigma_1$ behavior persists in important and widely used regimes, including
factorized initial conditions, weak coupling or short-time limits, and restricted
intervention classes for which discarded degrees of freedom never become causally
relevant. The present result clarifies the boundaries of these regimes rather than
challenging their empirical success.
\end{remark}

\subsection{Minimal unfaithful example with explicit unitary dynamics}

Consider a system $S$ coupled to an environment $E$, with a joint unitary evolution
$U_{SE}(\Delta t)$ representing free evolution over a duration $\Delta t$. Let the reduction
$\Pi$ retain only the system state, so $\Pi(h)$ corresponds to $\rho_S(t_0)$ at an operational
cut time $t_0$.

Define two admissible histories $h_1,h_2 \in \mathcal{H}$ whose joint states at $t_0$ are
pure but correlated in distinct ways:
\begin{align}
|\psi_1\rangle_{SE}(t_0) &= \sum_i \alpha_i \, |s_i\rangle_S \, |e_i^{(1)}\rangle_E, \\
|\psi_2\rangle_{SE}(t_0) &= \sum_i \alpha_i \, |s_i\rangle_S \, |e_i^{(2)}\rangle_E,
\end{align}
where $\{|s_i\rangle\}$ is an orthonormal basis on $S$ and the environment families
$\{|e_i^{(1)}\rangle\}$ and $\{|e_i^{(2)}\rangle\}$ differ (capturing distinct correlations).

Assume both histories reduce to the same system density operator:
\begin{equation}
\rho_S(t_0)
=
\mathrm{Tr}_E\!\left[ |\psi_1\rangle\langle\psi_1| \right]
=
\mathrm{Tr}_E\!\left[ |\psi_2\rangle\langle\psi_2| \right].
\end{equation}
In particular, in the basis $\{|s_i\rangle\}$ this reduced state has the operator form
\begin{equation}
\rho_S(t_0) = \sum_i |\alpha_i|^2 \, |s_i\rangle\langle s_i|,
\end{equation}
which makes explicit that $\rho_S(t_0)$ depends on the $\alpha_i$ but not on the environment
vectors $|e_i^{(k)}\rangle$. Thus $\Pi(h_1)=\Pi(h_2)=\rho_S(t_0)$, so the reduction is degenerate.

Let the intervention class $\mathcal{I}$ include the passive intervention $I_{\Delta t}$
corresponding to free evolution by $U_{SE}(\Delta t)$. The evolved joint states are
\begin{equation}
|\psi_k\rangle_{SE}(t_0+\Delta t)
=
U_{SE}(\Delta t)\,|\psi_k\rangle_{SE}(t_0),
\quad k \in \{1,2\}.
\end{equation}
The corresponding reduced states are
\begin{equation}
\rho_S^{(k)}(t_0+\Delta t)
=
\mathrm{Tr}_E\!\left[ |\psi_k\rangle\langle\psi_k| (t_0+\Delta t) \right].
\end{equation}
For generic interaction unitaries $U_{SE}(\Delta t)$ and distinct correlation structure
in the environment families, one obtains
\begin{equation}
\rho_S^{(1)}(t_0+\Delta t) \neq \rho_S^{(2)}(t_0+\Delta t),
\end{equation}
because $\rho_S(t_0)$ does not encode which environment correlation pattern is present.

This is a \textbf{predictive} failure: an experimenter who prepares $\rho_S(t_0)$ and knows only
the $\Sigma_1$ description cannot predict the outcome of applying $I_{\Delta t}$. A $\Sigma_1$ framework
would prescribe a unique reduced dynamical map $\mathcal{E}_{\Delta t}$ such that
\begin{equation}
\rho_S(t_0+\Delta t) = \mathcal{E}_{\Delta t}\!\left[\rho_S(t_0)\right],
\end{equation}
but in the presence of the unfaithful cut the same input $\rho_S(t_0)$ yields multiple distinct
outputs depending on hidden correlation structure. No choice of $\mathcal{E}_{\Delta t}$ can be
correct for both histories simultaneously without additional structure beyond the $\Sigma_1$ object.

\subsection{Classical Analog in Stochastic Control and Causal Inference}
\label{sec:classical-analog}

This subsection provides a non-quantum instantiation of the unfaithful cut.
It introduces no new claims for classical control or causal inference; it
serves only to demonstrate that unfaithful cuts arise whenever intervention
semantics are imposed on degenerate reductions, independent of quantum
structure.

\paragraph{Hidden-state control model.}
Let $X_t \in \mathcal{X}$ denote an unobserved state, $Y_t \in \mathcal{Y}$
an observed quantity, and $U_t \in \mathcal{U}$ a controllable intervention.
Assume a controlled stochastic evolution
\begin{equation}
  X_{t+1} = f(X_t, U_t, \xi_t),
  \qquad
  Y_t = g(X_t),
  \label{eq:classical-hidden-state}
\end{equation}
where $\{\xi_t\}$ is exogenous noise. A \emph{history} $h \in \mathcal{H}$
may be taken as a realization of $(X_0, Y_0, U_0, X_1, Y_1, U_1, \dots)$,
or more abstractly as an element of whatever admissible trajectory space
the model induces.

Define a reduction $\Pi:\mathcal{H}\to\mathcal{R}$ by retaining only the
current observation at a designated reduction event:
\begin{equation}
  \Pi(h) \coloneqq Y_t.
  \label{eq:classical-reduction}
\end{equation}
This is the classical analogue of working with an instantaneous reduced
descriptor and treating it as sufficient.

\paragraph{Interventions and continuation sets.}
Let the intervention class $\mathcal{I}$ include one-step control actions
$I_u$ that set $U_t=u$ at the next step and otherwise act only through the
retained degrees of freedom. Using the intervention-induced relation
$h \sim_{I} h'$ introduced in Section~\ref{sec:definitions}, define
continuation outcomes in the reduced space by
\begin{equation}
  \mathrm{Cont}(y;\mathcal{I})
  \;\coloneqq\;
  \Bigl\{ (I, y') \,:\,
  \exists\, h,h' \in \Pi^{-1}(y)\ \text{with}\ h \sim_{I} h'
  \ \wedge\ y'=\Pi(h') \Bigr\}.
  \label{eq:classical-cont}
\end{equation}
Faithfulness requires that if $\Pi(h_1)=\Pi(h_2)=y$, then
$\mathrm{Cont}(\Pi(h_1);\mathcal{I})=\mathrm{Cont}(\Pi(h_2);\mathcal{I})$.

\paragraph{An unfaithful cut from hidden confounding.}
Suppose there exist two admissible histories $h_1,h_2\in\mathcal{H}$ such that
\begin{equation}
  \Pi(h_1)=\Pi(h_2)=y,
  \qquad
  X_t^{(1)} \neq X_t^{(2)}.
  \label{eq:classical-same-observation}
\end{equation}
If the dynamics are correlation-sensitive in the sense that there exists a
control action $u$ for which
\begin{equation}
  \Pi(h_1') \neq \Pi(h_2')
  \quad \text{for some}\quad
  h_1 \sim_{I_u} h_1',\;\; h_2 \sim_{I_u} h_2',
  \label{eq:classical-diverge}
\end{equation}
then the reduction \eqref{eq:classical-reduction} is unfaithful: the same
reduced object $y$ supports distinct continuation structure under the same
intervention class $\mathcal{I}$.

This failure is the control-theoretic analogue of non-identifiability in causal
models: distinct latent states compatible with the same observation yield
different outcomes under identical do-interventions, as in Simpson's paradox driven by hidden confounding.

This is a predictive failure. Conditioning only on $Y_t=y$ does not suffice
to determine the outcome of applying $U_t=u$.

\paragraph{Concrete two-latent example.}
For intuition, take $\mathcal{X}=\{A,B\}$ and $\mathcal{Y}=\{0,1\}$ with a
degenerate observation map $g(A)=g(B)=0$. Let $\mathcal{U}=\{0,1\}$ and define
\begin{equation}
  \mathbb{P}(Y_{t+1}=1 \mid X_t=A, U_t=1)=1,
  \qquad
  \mathbb{P}(Y_{t+1}=1 \mid X_t=B, U_t=1)=0,
  \label{eq:classical-two-latent}
\end{equation}
(with any convenient specification for $U_t=0$). Then $Y_t=0$ is compatible
with both $X_t=A$ and $X_t=B$, but applying the same intervention $U_t=1$
yields deterministically distinct outcomes for $Y_{t+1}$.

\paragraph{$\Sigma_2$ resolution (classical form).}
As in the quantum case, the $\Sigma_2$ description replaces a point-valued
reduction by a fiber-aware object. The fiber at $y$ is
\begin{equation}
  \mathcal{F}(y) \coloneqq \Pi^{-1}(y),
  \label{eq:classical-fiber}
\end{equation}
equipped with the intervention-indexed admissibility relations inherited
from $\sim_I$. Interventions generically induce set-valued updates, reflecting
the fact that different admissible completions of the same reduced descriptor
support different future continuations.

\subsection{Addressing ``just add the environment''}

A natural response to the preceding example is to enlarge the reduced description to include
the environment, replacing $\rho_S$ by $\rho_{SE}$. This move can restore faithfulness for
that particular choice of ``system,'' but it does not eliminate the phenomenon: it merely
pushes the operational cut outward. Any finite retained subsystem can be correlated with
additional external degrees of freedom, and if the intervention class includes free evolution
of the larger closed dynamics, the same argument reappears.

\begin{corollary}[Pushing the cut outward does not remove unfaithfulness]
\label{cor:push-cut}
Let $\Pi$ describe any finite retained subsystem. If histories in $\mathcal{H}$ admit
correlations between the retained subsystem and discarded degrees of freedom, and if
$\mathcal{I}$ includes passive free evolution of the total system, then $\Pi$ is unfaithful
unless those correlations are excluded (or rendered dynamically irrelevant) by assumption.
\end{corollary}

Corollary~\ref{cor:push-cut} clarifies the structural content of the unfaithful cut: the issue
is not that one picked the ``wrong'' subsystem, but that any many-to-one operational reduction
that discards potentially relevant correlations will fail under sufficiently rich intervention
classes.